\documentclass[10pt,a4paper]{amsart}
\usepackage[margin=19mm]{geometry}
\usepackage{amsmath,amssymb,mathtools}
\usepackage[scr=rsfs,cal=euler]{mathalfa}
\usepackage{xfrac}
\usepackage[unicode,hidelinks,bookmarks=false]{hyperref}
\newcommand{\ie}{i.e.\ }
\newcommand{\cf}{cf.\ }
\newcommand{\resp}{respectively\ }
\newcommand{\Subsection}{\subsection}
\providecommand{\nobreakdash}{\nobreak\hbox{-}}
\setlength{\emergencystretch}{2em}
\multlinegap0pt
\usepackage{bm}
\usepackage{bbm}
\usepackage{dsfont}
\usepackage{xspace}
\usepackage{cancel}
\usepackage{mathtools}
\usepackage{cleveref}
\usepackage{amsthm}
\makeatletter
\@ifundefined{theorem}{\newtheorem{theorem}{Theorem}}{}
\@ifundefined{proposition}{\newtheorem{proposition}[theorem]{Proposition}}{}
\makeatother
\title[Densities for Elliptic Curves over Global Function Fields]{Densities for Elliptic Curves over Global Function Fields}
\author{Andrew Yao}
\date{Source: arXiv:2301.11437v2 (CC BY 4.0)}
\begin{document}
\maketitle

\noindent\emph{Source and licence.} Densities for Elliptic Curves over Global Function Fields. Version arXiv:2301.11437v2 (CC BY 4.0), distributed under the Creative Commons Attribution 4.0 International licence, as recorded in the arXiv licence metadata for this exact version and shown on the abstract page https://arxiv.org/abs/2301.11437v2. Full rights notice: licenses/arxiv-2301.11437-CC-BY-4.0.txt.\par\smallskip

\noindent\textit{Editorial scope.} Counted unit: the Theorem whose source label is \texttt{prop:unique2} in the article (source0.tex lines 714--717, source label \texttt{prop:unique2}), delivered together with the article's own complete original proof of it (source0.tex lines 719--774, one \texttt{proof} environment of 56 lines). No mathematics is written, repaired or re-derived: statement and proof are the article's own text line for line. Required context retained uncounted: translation (lines 165--168); prop:multiterations (lines 173--184). Every definition, display and label of those lines is retained unchanged. Omitted from this package and not redistributed: 1--164, 169--172, 185--713, 718--718, 775--1211. Nothing inside a retained excerpt is omitted or altered: every retained body is delivered exactly as the article's lines are written, so no omission receipt of any kind accompanies this package. Citations: the retained excerpts carry no citation command at all, so no bibliography entry of the article is transcribed and no reference is redistributed. Reference adaptation: none was needed and none was made; every reference inside the delivered unit resolves inside it, so no unresolved reference or citation is produced. Printed slips kept literal: no printed word, symbol, hypothesis or number of the counted statement or of its proof was altered. Layout and class: the article's own class is \texttt{amsart} with packages including geometry, babel, amsrefs, amsmath, amssymb, amsthm, amscd, bbm, esvect, mathrsfs, asymptote, hyperref, cleveref, pgf; the wrapper is the lane's pinned \texttt{amsart} editorial wrapper at 10pt on A4, which restates the theorem-like environments the retained text needs and adds the article's own macro definitions that the retained text needs (none), with extra wrapper packages (bm, bbm, dsfont, xspace, cancel, mathtools, cleveref). The delivered numbering is the unit's own. Assets: no retained span contains a figure environment, a graphics-inclusion command, a PDF-page inclusion, a table or a TikZ picture; no image file is copied into this package, the package declares no asset dependency and no image file is redistributed with it. Licence: version arXiv:2301.11437v2 of this article is distributed under the Creative Commons Attribution 4.0 International licence, as recorded in the arXiv licence metadata for this exact version and shown on the abstract page https://arxiv.org/abs/2301.11437v2. A full rights notice is delivered at licenses/arxiv-2301.11437-CC-BY-4.0.txt. Characters: every retained excerpt is pure ASCII, so the delivered engine, XeLaTeX with its default Latin Modern text font, reports no missing character.
\begin{proposition}
\label{translation}
Let $E$ and $F$ be elliptic curves over $K_P$ that have equations with coefficients in $R_P$. Assume that $E$ and $F$ are isomorphic and satisfy $v_P(\Delta(E))=v_P(\Delta(F))$. Then, there exists $l,m,n,u\in R_P$ such that $v_P(u)=0$ and the equation of $F$ can be obtained from the equation of $E$ by first replacing $x$ with $u^2x+n$ and $y$ with $u^3y+lu^2x+m$ and then dividing by $u^6$.
\end{proposition}

\begin{proposition}
\label{prop:multiterations}
     Let $k$ be a nonnegative integer. Suppose $E$ is an elliptic curve over $K_P$ with equation
     \[
     E: y^2+a_1xy+a_3y=x^3+a_2x^2+a_4x+a_6
     \]
     and assume that $a_1,a_2,a_3,a_4,a_6\in R_P$. For $l,m,n\in K_P$, let $E'(l,m,n)$ be the elliptic curve that is $E$ with $x$ replaced by $x+n$ and $y$ replaced by $y+lx+m$. Then, $N_P(E)\geq k$ if and only if there exists $l,m,n\in R_P$ such that if $E'(l,m,n)$ has equation
     \[
     E'(l,m,n):y^2+a_1'xy+a_3'y=x^3+a_2'x^2+a_4'x+a_6',
     \]
     where $a_i'\in \pi_P^{ki}R_P$ for $i\in\{1,2,3,4,6\}$.
\end{proposition}

\begin{theorem}
\label{prop:unique2}
Let $E$ be an elliptic curve in $G_P^{(3)}$. Then, $E\in S_k$ if and only if a unique pair $(l,m)\in L_{P,k}\times L_{P,3k}$ exists such that $(l,m,l^2+a_1l)\in A_k(E)$.
\end{theorem}

\begin{proof}
Suppose a unique pair $(l,m)$ satisfying the conditions exists. Because $A_k(E)$ is nonempty, $E\in S_k$ from \Cref{prop:multiterations}.

Assume $E\in S_k$. Then, using \Cref{prop:multiterations}, $A_k(E)$ is nonempty. Let the equation of $E$ be $E:y^2+a_1xy+a_3y=x^3+a_4x+a_6$ for $a_1,a_3,a_4,a_6\in R_P$.

From replacing $y$ with $y+l'x$ for $l'\in R_P$, if $(l,m,n)\in A_k(E)$, $(l+l'\pi_P^{k}, m,n)\in A_k(E)$. Therefore, there exist $l\in L_{P,k}$ and $m,n\in R_P$ such that $(l,m,n)\in A_k(E)$. Moreover, if $(l,m,n)\in A_k(E)$, $l^2+a_1l+n\equiv 0\pmod{\pi_P^{2k}}$. With this, from replacing $x$ with $x+\frac{l^2+a_1l+n}{\pi_P^{2k}}$, if $(l,m,n)\in A_k(E)$, $(l,m+l(l^2+a_1l+n),l^2+a_1l)\in A_k(E)$. Therefore, there exist $l\in L_{P,k}$ and $m\in R_P$ such that $(l,m,l^2+a_1l)\in A_k(E)$. Next, from replacing $y$ with $y+m'$ for $m'\in R_P$, there exists $l\in L_{P,k}$ and $m\in L_{P,3k}$ such that $(l,m,l^2+a_1l)\in A_k(E)$.

Next, we prove that $(l,m)$ is unique. Assume that $(l_1,m_1), (l_2,m_2)\in L_{P,k}\times L_{P,3k}$ and $(l_1,m_1,l_1^2+a_1l_1), (l_2,m_2,l_2^2+a_1l_2) \in A_k(E)$. We prove that $(l_1,m_1)=(l_2,m_2)$.

Let $F$ be the curve 
\[
F:y^2+\frac{a_1}{\pi_P^k}xy+\frac{a_3}{\pi_P^{3k}}y=x^3+\frac{a_4}{\pi_P^k}x+\frac{a_6}{\pi_P^{6k}}.
\]
For $1\leq i\leq 2$, let $F_i$ be $F$ with $x$ replaced by $x+\frac{l_i^2+a_1l_i}{\pi_P^{2k}}$ and $y$ replaced by $y+\frac{l_i}{\pi_P^{k}}x+\frac{m_i}{\pi_P^{3k}}$. Note that $F_i\in G_P^{(3)}$ because $(l_i, m_i, l_i^2+a_il_i)\in A_k(E)$ for $1\leq i\leq 2$. From this, $a_1\equiv 0\pmod{\pi_P^k}$. 

Suppose $a_1\not=0$. We have that $F_1$ and $F_2$ are isomorphic and $v_P(\Delta(F_1))=v_P(\Delta(F_2))$. Then, using \Cref{translation}, let $\tau$ be a translation from the equation of $F_1$ to the equation of $F_2$ that replaces $x$ with $u^2x+n'$ and $y$ with $u^3y+l'u^2x+m'$, where $u,l',m',n'\in R_P$ and $v_P(u)=0$. 

The coefficient of $xy$ after $\tau$ is applied to the equation of $F_1$ is $\frac{a_1}{u\pi_P^k}$. However, the coefficient of $xy$ in $F_2$ is $\frac{a_1}{\pi_P^k}$. Therefore, $u=1$ and $a_1\equiv 0\pmod{\pi_P^k}$. 

Next, the coefficient of $y$ after $\tau$ is applied to the equation of $F_1$ \[
\frac{a_1l_1^2+a_1^2l_1+a_3+\pi_P^{2k}a_1n'}{\pi_P^{3k}}.
\]
However, the coefficient of $y$ in $F_2$ is
\[
\frac{a_1l_2^2+a_1^2l_2+a_3}{\pi_P^{3k}}.
\]
Therefore, 
\[
l_1^2+a_1l_1+\pi_P^{2k}n'=l_2^2+a_1l_2.
\]
Because $a_1\equiv 0\pmod{\pi_P^k}$, we have that $l_1\equiv l_2\pmod{\pi_P^k}$. Therefore, $l_1=l_2$. From this, $n'=0$.

The coefficient of $x^2$ after $\tau$ is applied to the equation of $F_1$ is
\[
n'+(l')^2+\frac{a_1l'}{\pi_P^k}.
\]
This equals the coefficient of $x^2$ in $F_2$, which is $0$. Because $n'=0$, we have that $l'=0$ or $l'=\frac{a_1}{\pi_P^k}$.

From setting the coefficient of $x$ after $\tau$ is applied to the equation of $F_1$ equal to the coefficient of $x$ in $F_2$,
\[
\frac{a_1}{\pi_P^k}\cdot\left(\frac{m_1}{\pi_P^{3k}}+m'\right)+\frac{a_1(l_1^2+a_1l_1)+a_3}{\pi_P^{3k}}\cdot l'=\frac{a_1}{\pi_P^k}\cdot\frac{m_2}{\pi_P^{3k}}.
\]
Suppose $l'=0$. Then $\frac{m_1}{\pi_P^{3k}}+m'=\frac{m_2}{\pi_P^{3k}}$. It follows that $m_1\equiv m_2\pmod{\pi_P^{3k}}$ and $m_1=m_2$. Suppose $l'=\frac{a_1}{\pi_P^k}$. We have that
\[
\frac{m_1}{\pi_P^{3k}}+m'+\frac{a_1(l_1^2+a_1l_1)+a_3}{\pi_P^{3k}}=\frac{m_2}{\pi_P^{3k}}.
\]
However, using that the coefficient of $y$ in $F_2$ is an element of $R_P$,
\[
a_1(l_1^2+a_1l_1)+a_3\equiv a_1(l_2^2+a_1l_2)+a_3\equiv 0\pmod{\pi_P^{3k}}.
\]
Therefore, $m_1\equiv m_2\pmod{\pi_P^{3k}}$ and $m_1=m_2$.

Assume $a_1=0$. From the coefficient of $y$ in $F_2$, we have that $a_3\equiv 0\pmod{\pi_P^{3k}}$. Also, from the coefficients of $x$ in $F_1$ and $F_2$,
$l_1^4+a_3l_1\equiv l_2^4+a_3l_2\pmod{\pi_P^{4k}}$.
This gives that $l_1=l_2$. Afterwards, from the constant terms of $F_1$ and $F_2$, $m_1^2+a_3m_1\equiv m_2^2+a_3m_2\pmod{\pi_P^{6k}}$. From this, we obtain that $m_1=m_2$.
\end{proof}

\end{document}
